\section{Appendix: Scope and Qualifications}
\label{sec:scope}

{\small
\textbf{Edition, formulary and occurrence.} The study treats the Mass
\latin{Dominica vigesima post Pentecosten}, II classis, of the \work{Missale
Romanum}, editio typica, Typis Polyglottis Vaticanis 1962, pp.~412--413,
nos.~1689--1698, in the universal 1962 calendar and no particular calendar. Its
occurrence on 11~October 2026, its precedence over the feast of the Maternity
of the Blessed Virgin Mary, the commemoration of that feast in low and
conventual Masses and its omission in a sung Mass that is not conventual, the
colour, the Credo and the Trinity Preface rest on the Missal's own calendar and
general rubrics (RG 16, 91, 108, 111, 115, 127~b; RGMR 434, 437~a, 475~a, 483,
494~b), read on the page images of the typical edition. The general rubrics tie
the Gloria to the Te Deum of the Office (RGMR 431~a); the Breviary rubric behind
it, which the typical edition of the Missal does not print, was not read. The
commemorated orations are given from the Maternity's formulary, pp.~685--686;
nothing else of that formulary enters the Sunday Mass. No diocesan, titular,
dedication or institute calendar, and no local external solemnity, is applied.

\textbf{Text control.} Every Latin form of the ten proper elements and of the
three commemorated orations was read on page images of the CMAA facsimile of
the typical edition and collated with the Pustet Missal of 1862, in which every
element's wording already stands; the collation, the conclusion forms and the
comparison with the Clementine Vulgate are recorded in
\texttt{propers/verified.md}. The Introit's antiphon is a cento of words from
Dan 3:29--43 that no witness read carries in the Missal's form, and the Latin
source of the Gradual's \latin{illis escam}, the Alleluia's \latin{psallam tibi,
gloria mea} and the Communion's added \latin{Domine} was not found in the
psalter texts of the Vulgate, the Greek Psalter and the Latin Fathers read; the
Offertory's \latin{dum recordaremur tui, Sion} stands in Cassiodorus's psalm
text. No Roman or Old Latin psalter and no chant manuscript was consulted.

\textbf{English.} The scriptural English is the Douay--Rheims in Challoner's
revision at the canonical verses. Where the Missal departs from the Vulgate (the
Introit's antiphon, the Epistle's opening, the Gradual's \latin{illis}, the
Alleluia's close, the Gospel's opening, the Offertory's \latin{tui}, and the
Communion's \latin{Memento}, \latin{Domine} and shortened close) that English
does not render the Missal's Latin, and no English has been supplied in its
place. The English of the Introit's antiphon and of the Sunday's three orations
is the anonymous translation printed by Cummiskey in Philadelphia in 1861: the
orations in the library's tracked transcription, checked against the scan's page
images, and the antiphon read on the page image of printed p.~448. The English
of the commemorated orations is the 1945 Benziger revision of Lasance's
\work{The New Roman Missal}. None of these is an approved liturgical translation
of the 1962 Missal. The Trinity Preface is cited and not reproduced. The
English of St Augustine's tractate on St John and his expositions of the
psalms, of St John Chrysostom on St John and on Ephesians, and of St Cyril of
Alexandria is that of the nineteenth-century Nicene and Post-Nicene Fathers and
Library of the Fathers. Where the Latin of St Gregory, St Jerome, St Hilary and
St Augustine is printed, the English beside it is a gloss of that Latin and not
a translation from a named edition.

\textbf{Reception.} Each appointed passage was searched in the Latin and Greek
Fathers and the later Doctors, saints and liturgical commentators available to
this production, and the witnesses read are those named in the References at
their exact loci. St Augustine's expositions of Pss~118, 136 and 144 were read in
English and, where the Latin is quoted, on the page images of Migne's
\work{Patrologia Latina}; St Hilary's tractates on Pss~118 and 136 and St
Jerome's commentary on Ephesians on Migne's page images; St Gregory's homily in
the Latin Wikisource transcription of Migne's text, read beside the page images
of the printed columns. St Thomas Aquinas, the medieval liturgical commentators
and Cassiodorus were read in text layers or in a restricted scan and are
described, not quoted; so is St John Chrysostom's Greek exposition of
Pss~136 and 144, and St Anthony of Padua's sermon, read only in a copyrighted
English translation. Rupert of Deutz, Honorius Augustodunensis, Sicard of
Cremona and William Durandus are cited for what they say of the Introit,
Epistle, Gradual, Offertory and Communion, which the Mass they expounded shared
with this formulary although it read another Gospel; Rupert and Durandus are
also cited on this formulary's Gospel at the chapters where their books
expound it. Bl.\ Ildefonso
Schuster and the continuation of \work{The Liturgical Year} comment on this
formulary as the 1962 Missal gives it. Among the witnesses not reached are
Origen's commentary on St John, St Ambrose on Ps~118, Theodoret and St
Athanasius on the psalms, Theophylact on Ephesians, the Old Latin and Roman
psalters, the Breviary's homily for this Sunday, and St Alphonsus Liguori's
Sunday sermons.

\textbf{Chronology.} The Date cells in Scriptural Date and Location come from
the Scripture chronology corpus, through the record generated for this
formulary: its default profile's answer for every element and, for the
Epistle, a declared comparison under its critical profile. For the four psalm
elements it gives the modern critical bound for the whole Psalter, from the
introduction to the Psalms in the \work{New American Bible, Revised Edition}
(summarised, not quoted), and, for Pss~107, 118 and 144, David's reign in the
usual chronology of the \work{Catholic Encyclopedia}'s ``David'', disputed
because the article reports a later dating beside it, as the reference point of
the attribution and not as a date of any psalm. For Ps~136 it gives a setting,
the Babylonian captivity, with the seventy years as a duration and the
destruction of Jerusalem as historical background. For the Introit it holds no
date common to Daniel and the psalm verse; its explanatory row gives each
passage its own answer. For the Gospel it holds a derived event year, from
A.~J. Maas's ``Jesus Christ'' in the encyclopedia, and two composition dates,
both disputed; for the Epistle, two traditional dates, both disputed, and the
comparison with the hypothesis of a later disciple, which does not enter the
default answer. The record and its sources are audited in
\texttt{research/scope.md}.

\textbf{Rights.} The Latin of the formulary and of the commemorated orations is
published on the basis of its public-domain antecedent, the Pustet Missal of
1862 (17 U.S.C. 103(b)). The Douay--Rheims, the 1861 Cummiskey English, Migne's
Latin and Greek texts, the Nicene and Post-Nicene Fathers, the Library of the
Fathers' Cyril, the Venice 1745 and Liège 1857 Aquinas, O'Sullivan's English
of Bellarmine, the 1883 English of the continuation of \work{The Liturgical
Year} and the 1927 English of Schuster's \work{Sacramentary} are in the public
domain in the United States; the 1945 Lasance English rests on the nonrenewal
recorded in the source library. St Gregory's Latin is quoted from the Latin
Wikisource transcription of the \work{Homiliarum in Evangelia}, licensed under
Creative Commons Attribution-ShareAlike 3.0, with attribution to its
contributors at \url{https://la.wikisource.org/wiki/Homiliarum_in_Evangelia}.
Cassiodorus in the \work{Corpus Christianorum} edition and St Anthony in
Spilsbury's English are described, not reproduced; the New American Bible
introductions are summarised, not reproduced.

\textbf{Records.} The collation of the appointed texts is recorded in
\texttt{propers/verified.md}, and the occurrence and rubrics in
\texttt{research/context.md}. The search, loci, disagreements and negative
results are in \texttt{research/scope.md}, and the reasoning behind the three
readings in \texttt{research/interpretations.md}.
The research records were reviewed independently before this study was written.
\par}

\phantomsection
\section*{References}
% A starred heading sets no mark, so the Appendix that ends above would go on
% heading the pages of the bibliography. The heading stays starred, because
% check-content-preflight finds the References section by that exact form; the
% mark is set explicitly through the leaf-local \runninghead, which the web
% converter's macro whitelist carries.
\runninghead{References}
\addcontentsline{toc}{section}{References}
\label{sec:references}

{\footnotesize
\subsection*{Liturgical books and Bibles}
\begin{itemize}
\item \work{Missale Romanum ex decreto Sacrosancti Concilii Tridentini restitutum, Summorum Pontificum cura recognitum}, editio typica (Typis Polyglottis Vaticanis, 1962), CMAA facsimile, \url{https://media.churchmusicassociation.org/pdf/missale62.pdf}: \work{Rubricae generales} 16, 91, 108, 111, 115, 127; \work{Rubricae generales Missalis} 427, 431, 434, 437, 475, 483, 494; \work{Ordo Missae} no.~1033; pp.~412--413 (nos.~1689--1698); pp.~685--686 (nos.~3847, 3855, 3857).
\item \work{Missale Romanum} (Ratisbon: Pustet, 1862), pp.~352--354, \latin{Dominica XX. post Pentecosten}; appendix pp.~[171]--[172], \latin{Festum Maternitatis B.\ Mariae Virginis}.
\item \work{The Roman Missal, Translated into the English Language for the Use of the Laity} (Philadelphia: Eugene Cummiskey, 1861), pp.~448--450 (XX Sunday after Pentecost, Introit, Collect, Secret, Postcommunion).
\item F.~X. Lasance, \work{The New Roman Missal}, revised edition (New York: Benziger Brothers, 1945), pp.~1233, 1235--1236 (Maternity of the Blessed Virgin Mary, Collect, Secret, Postcommunion).
\item The Holy Bible, Douay--Rheims version, revised by Bishop Richard Challoner (Project Gutenberg text): Deut 32:15; 2 Mac 15:14; Ps 56:1, 5, 8--9; 59:1--2; 107:1--4; 118:1--3, 19, 49--52; 136:1--9; 144:1, 13--18; Prov 4:18; Dan 3:1, 20, 24--51; Jn 4:39--54; Eph 1:1; 3:1; 4:1; 5:8--33; 6:20.
\item \work{Biblia Sacra Vulgatae editionis} (Clementine Vulgate), eBible.org text: Ps 107:2--4; 118:1, 49--50; 136:1; 144:15--16; Dan 3:29--43; Jn 4:46--53; Eph 5:15--21.
\end{itemize}

\subsection*{Fathers, Doctors and saints}
\begin{itemize}
\item St Augustine, \work{Tractates on the Gospel of John} 16, 3--7, trans. in Nicene and Post-Nicene Fathers, 1st ser., vol.~7 (CCEL text).
\item St Augustine, \work{Enarrationes in Psalmos}, trans. in Nicene and Post-Nicene Fathers, 1st ser., vol.~8 (CCEL text): on Ps~56:8; 118, sermons 1 and 15 (on vv.~1--3, 49--50); 136, on vv.~1--9; 144, on vv.~14--17. Latin in Migne, PL 37 (Paris, 1845), cols.~1502 (118, sermo 1, 1), 1761 and 1774 (136), 1882 (144).
\item St Gregory the Great, \work{Homiliae in Evangelia} 4, 3 and 28, 1--3, Latin Wikisource transcription of Migne's text, CC BY-SA 3.0; Migne, PL 76, cols.~1210--1213.
\item St Jerome, \work{Commentarii in Danielem}, preface and on Dan 3:23--49, in Migne, PL 25 (Latin Wikisource transcription), cols.~639--642; \work{Commentarii in Epistolam ad Ephesios} III, on Eph 5:16--20, in Migne, PL 26 (Paris, 1845), cols.~527--530.
\item St Hilary of Poitiers, \work{Tractatus super Psalmos}, in Ps.~118, \latin{Aleph} 1--3 and \latin{Zain}; in Ps.~136; in Ps.~144, in Migne, PL 9 (Paris, 1844), cols.~503--505, 548--550, 777--784, 863--864.
\item St John Chrysostom, \work{Homilies on St John} 35, Nicene and Post-Nicene Fathers, 1st ser., vol.~14 (New York, 1889), pp.~123--125; \work{Homilies on Ephesians} 19, Nicene and Post-Nicene Fathers, 1st ser., vol.~13 (CCEL text); \work{Expositio in Psalmos}, in Ps.~136 and 144, Migne, PG 55, cols.~405--407 and 470--471 (text layer; described, not quoted).
\item St Cyril of Alexandria, \work{Commentary on the Gospel according to S.~John} II.5, trans. P.~E. Pusey, Library of the Fathers (Oxford, 1874), pp.~232--236.
\item St Thomas Aquinas, \work{Super Evangelium S.~Ioannis lectura} c.~4, lect.~7, in \work{Opera}, editio altera Veneta, vol.~3 (Venice, 1745), pp.~521--524; \work{In Epistolam ad Ephesios} c.~5, lect.~6, in \work{In omnes D.~Pauli Apostoli epistolas commentaria}, vol.~2 (Liège: H.~Dessain, 1857), pp.~344--345 (text layers; described, not quoted).
\item St Robert Bellarmine, \work{A Commentary on the Book of Psalms}, trans. J.~O'Sullivan, read in the eCatholic2000 transcription of that translation: on Pss~56:7--8; 118:1--3, 49; 136; 144:15--16.
\item St Anthony of Padua, Sunday sermon on Jn 4:46--53, in \work{Sermons for Sundays and Festivals}, trans. S.~R.~P. Spilsbury (described, not quoted).
\item Cassiodorus, \work{Expositio Psalmorum}, on Pss~107:2--4; 118, preface and vv.~1, 49--50; 136; 144:15--16, ed. M.~Adriaen, CCSL 98 (Turnhout, 1958), and in Migne, PL 70 (described, not quoted).
\end{itemize}

\subsection*{Liturgical commentators}
\begin{itemize}
\item Bl.\ Ildefonso Schuster, \work{The Sacramentary (Liber Sacramentorum)}, vol.~3 (London: Burns Oates \& Washbourne, 1927), pp.~174--178.
\item The continuation of Prosper Guéranger, \work{The Liturgical Year} (Dom Lucien Fromage), \work{The Time after Pentecost}, vol.~2 (series vol.~11), trans. Laurence Shepherd (Dublin: James Duffy and Sons, 1883), pp.~439, 447, 451--452.
\item Rupert of Deutz, \work{De divinis officiis} XII.20--21, in Migne, PL 170 (text layer).
\item Honorius Augustodunensis, \work{Gemma animae} IV.88--89, in Migne, PL 172 (text layer).
\item Sicard of Cremona, \work{Mitrale} VIII.20, in Migne, PL 213 (text layer).
\item William Durandus, \work{Rationale divinorum officiorum} VI.137--138, vol.~2 (Lyon, 1612; text layer).
\end{itemize}

\subsection*{Chronology sources}
\begin{itemize}
\item \work{The Catholic Encyclopedia} (New York, 1907--1912): ``Captivities of the Israelites'' (vol.~3, 1908); ``Book of Daniel'' and ``David'' (vol.~4, 1908); ``Epistle to the Ephesians'' (vol.~5, 1909); ``Jesus Christ'', ``The Gospel of Saint John'', ``Third and Fourth Books of Kings'' and ``Chronology of the Kings'' (vol.~8, 1910); ``St.\ Paul'' (vol.~11, 1911); ``The New Testament'' and ``Temple of Jerusalem'' (vol.~14, 1912).
\item G.~L. Haydock, \work{The Douay--Rheims Bible with Haydock's Commentary}, notes on Pss~70 and 136.
\item United States Conference of Catholic Bishops, \work{New American Bible, Revised Edition}, introductions to the Psalms and to Ephesians, on their dating; summarised, not reproduced.
\end{itemize}
\par}
